\section{The agent contract and its governance}\label{sec:contract}

An agent is an external program in any language that reads a \texttt{MarketObservation} as JSON and replies with a \texttt{Decision} as JSON. The harness sends an observation, the agent returns a decision, repeat. Everything else in this paper exists to make that loop trustworthy.

\subsection{Authoritative artifacts}\label{sec:contract-artifacts}

The contract has four authoritative artifacts, all committed. \texttt{CONTRACT\_VERSION}, currently \texttt{1.0}, is the single source of truth for the wire version and lives in the contract module of the core crate. Two JSON Schemas, draft 2020-12, define \texttt{MarketObservation} and \texttt{Decision} for non-Rust implementers; \cref{sec:spaces} enumerates their fields. A conformance kit of fixture documents, covering the normal multi-symbol case, a single symbol, an empty-portfolio start, signed-weight shorting and a legacy decision shape, is exercised by a committed conformance test. And a governance document states the evolution rules in writing. The schema closes the JSON shape but cannot know the symbols in the current observation, so one canonical runtime parser additionally enforces episode semantics. The MCP, training and local-agent surfaces all call it; a semantic fault returns without advancing the environment.

\subsection{Additive-only evolution}\label{sec:contract-evolution}

The contract evolves additively. The only backwards-compatible change is a new field that is optional with a default, so an agent written against an older version keeps parsing newer observations and the harness keeps parsing older decisions. The optional fields are absent from the schemas' \texttt{required} lists by construction. Removing, renaming or retyping a field, making an optional field required, and removing or renaming an action variant are all defined as breaking and forbidden on the v1 surface. A breaking change never mutates the v1 types in place; it ships a parallel namespace with its own schemas and its own major version, and the two run side by side through a published deprecation window. Adding a new action variant is itself classified as breaking, because consumers that exhaustively match would misparse it, so it also goes through the parallel namespace rather than an additive bump.

\subsection{Version semantics}\label{sec:contract-versions}

\texttt{CONTRACT\_VERSION} versions the wire shape only and is deliberately decoupled from the crate and package versions, which move on their own release cadence. A major bump is a breaking change shipped as a parallel namespace. A minor bump is reserved for a batch of additive fields significant enough to advertise; a single additive field needs no bump at all, because existing agents are unaffected by definition. A package release never, on its own, bumps the contract version.

\subsection{A transport fault is a failed cell, not a hold}\label{sec:transport-gate}

An external agent speaks over a wire, and a wire fails. The v1.0 \texttt{Decision} carries no error variant, and the external transports cannot signal one through the \texttt{Agent} trait they implement, so a timeout, a dropped connection or an unparseable reply degrades to an empty order list while the fault itself is recorded into a rolling \texttt{TransportHealth}. \Cref{sec:spaces} fixes what an empty order list means: it is a hold, and against a stateful engine a hold is not inaction. The position persists. An agent whose wire wedges on the first bar therefore rides its opening position to the end of the window, the run completes without raising, and the return series it emits is the return series of a deliberately conservative agent. No downstream gate separates the two, because the deflation, reliability and process legs all consume returns. A committed characterization test states the failure rather than describing it: a wedged agent run through the plain engine call yields a scoreable \texttt{Run}.

The health record mitigates this failure only when a scored path reads it. The crate therefore re-exports the diagnostics types, and \texttt{run\_backtest\_checked} is the explicit Rust entry point for an external agent: it runs the cell, reads the health afterwards, and returns a \texttt{CellOutcome} that is either \texttt{Scored(Run)} or \texttt{Failed(TransportFault)}. The local-field scheduler and artifact compiler enforce the same distinction in their own typed cell schema, so their faults cannot enter a scored artifact as flat returns. The low-level compatibility function \texttt{run\_backtest} remains public and degrades an external transport fault to the transport's empty-decision fallback; the type system does not force a wire consumer to choose the checked entry point. The guarantee therefore covers the checked and repository scoring paths, not every possible caller of the compatibility API. A typed failure carries the retryable transport-fault count, the agent-protocol fault count, whether the per-endpoint circuit breaker tripped, and the last decision-level error. Its accessor returns \texttt{None} on a failed cell rather than a default-valued \texttt{Run}, so a caller that chose the checked path cannot score a failure by accident. Tests preserve both sides of the boundary: a wedged agent through the compatibility call produces a scoreable run, whereas the checked call yields a failed cell; clean and malformed transports exercise the other outcomes.

The fix belongs with the determinism and replay guarantees rather than with operations, because it is the same kind of claim. Leak-freedom bounds what the agent could see and replay bounds what the trajectory could be; neither says anything about a return series the harness synthesized on the agent's behalf. Masking is the one failure mode that produces evidence a verifier cannot reject, since the doctored quantity originates inside the honest producer and replays byte-for-byte. A failure is evidence, and never a hold.

\subsection{Why the contract can be trusted without trusting the host}\label{sec:contract-trust}

The governance promise is social; the properties of \cref{sec:leakfree,sec:determinism,sec:replay} are technical, and together they remove the host from the trust equation for everything the byte-identity guarantee covers. That coverage has a boundary. The Rust core, its WebAssembly build and the scenario bytes returned through the Python bindings are pinned to the byte by committed FNV-1a goldens, as are the engine outputs behind F1 through F3. The native and Python suites assert the full set. The canonical and clustered scenarios are recorded once, in \texttt{crates/sharpearena/contract/attestation/scenario-goldens.json}, and three gates read that one file: the native suite, a \texttt{wasm-pack} run that executes the exported entry point inside a WebAssembly module under Node, and the npm package's own suite, which asserts the committed \texttt{.wasm} binary it ships rather than a host build of the same source. A generator change therefore cannot reach a release through one runtime alone. The Python evidence layer behind the remaining probes is deterministic in its seeds but sits outside the golden-hash path (\cref{sec:limitations}). Within the covered surface, a leaderboard entry can be recomputed from the frozen seed and submitted decisions on any tested platform. SharpeArena exposes the published scoring helpers used by its diagnostics, while the SharpeBench kernel independently defines board-level deflation, reliability, process gates and rank eligibility \citep{toca2026sharpebench,bailey2014deflated}.
